(1) 速度分别是 \(\dot X\) 与 \(\dot X+\dot q\)，故 \(L=M\dot X^2/2+m(\dot X+\dot q)^2/2-kq^2/2\)。坐标 \(X\) 不出现在 \(L\) 中，广义动量 \(P=\partial L/\partial\dot X=(M+m)\dot X+m\dot q\) 守恒。令质心 \(C=X+mq/(M+m)\)，则 \(P=(M+m)\dot C\)，这是解耦的关键，单独令车静止会违背总动量守恒。

(2) 把 \(\dot X=(P-m\dot q)/(M+m)\) 代入动力学，得到 \(\mu\ddot q+kq=0\)，其中约化质量 \(\mu=mM/(M+m)\)。记 \(\omega=\sqrt{k/\mu}\)，唯一满足初值的解是 \(q(t)=q_0\cos\omega t+(u_0/\omega)\sin\omega t\)，而 \(C(t)=C_0+Pt/(M+m)\)。相对振幅为 \(\sqrt{q_0^2+(u_0/\omega)^2}\)。

(3) 由 \(X=C-mq/(M+m)\)、\(x_m=C+Mq/(M+m)\) 可直接给出两者轨迹；其速度分别是 \(P/(M+m)-m\dot q/(M+m)\) 与 \(P/(M+m)+M\dot q/(M+m)\)。代入得到 \(M\dot X+ m\dot x_m=P\) 对所有时刻成立。若 \(P=0\)，两个物体始终反相移动，振幅比为 \(m/M\)，而非两个物体具有相同振幅。

(4) 对动能配方得到 \(K=P^2/[2(M+m)]+\mu\dot q^2/2\)，所以总机械能 \(E=P^2/[2(M+m)]+\mu\dot q^2/2+kq^2/2\) 严格守恒。第二部分可用 \(\dot q^2+\omega^2q^2=\omega^2q_0^2+u_0^2\) 核实不变。\(M\to\infty\) 且 \(P/M\to0\) 时，\(\mu\to m,\omega\to\sqrt{k/m}\)，滑车反冲消失，恢复固定弹簧一端的标准频率。
